\section{结论}

本文面向缺乏统一样本级故障真值的多通道中子电流历史档案，将研究目标从自动故障分类转为可追溯的诊断辅助：先将原始测量转化为工程师可复核的观察，再使后续技术解释能够回指至适当的资料依据。为此，本文构建了 measurement-to-evidence interface。测量侧以确定性脚本生成并保存 observation packet；紧凑摘要直接由稳定的结构化字段形成，作为证据侧的受控检索意图；语言模型只组织已获得的观察、资料段落和候选解释，不从原始波形自行生成测量事实，也不确认物理机理。

代表性 trace 表明，这一接口能够把长时、多通道记录压缩为可检查的观察链。A2 记录中的孤立零值、短时偏离、通道间时间关系和极值位置均可回指到相应的输入、参数与工具输出。跨通道重复的零值时间因而可被保留为测量完整性或共享仪表上下文的检索线索，而不会被直接提升为某一种已确认故障。该结果说明，在没有统一标签的条件下，观测的可回指性本身构成了后续工程复核的必要基础。

证据侧实验表明，语义相似度不能单独承担来源治理的角色。在 \texttt{nomic-embed-text} 字符分块扩库协议中，相关补充资料加入后，最高余弦相似度可以提高或保持稳定，而前 10 个结果中核心权威资料的占比却会下降。MiniLM encoder 对照中同样观察到这一风险，说明该现象不宜归因于单一向量表示模型。资料范围、query 表述、语言、encoder 和分块方式共同改变检索结果；尤其是，技术细节 query 并不必然提高前 10 个结果中核心资料的占比。与此同时，实验能够按扩库后的 CorePriority@10 筛出这一占比保持较高的候选组合，为检索协议的条件性选择提供依据。由此，补充资料不应被视为无关噪声，但若应用要求特定权威资料集合保持在检索结果前列，来源 metadata 与来源策略必须被明确纳入检索设计和评价。

检索单元的核查还表明，重叠分块下“找回同一段落”与“找回同一文档”回答不同问题。同文档检索表现高于严格同分块检索，说明相邻分块可能保留了所需内容；因此，分块边界、段落定位与来源可见性需要分别报告，而不能以单一检索分数替代。综合这些结果，本文的结论是：无标签工程时序的安全使用依赖于“先描述、后判断”的责任分工——确定性工具产生可回指观察，检索系统定位带来源信息的技术证据，语言模型组织受证据约束的候选解释，工程师保留机理判断和最终决策。后续若能接入脱敏的运行/维护背景并获得领域专家复核，可在该接口之上进一步检验机理假设、段落支持与工程有用性。
